	\vspace{12pt}
	\section{Vorwort}
		Die Karate Abteilung am Standort Reusrath ist Mitglied der \href{https://www.sglangenfeld.de/}{Sportgemeinschaft Langenfeld}, kurz \href{https://www.sglangenfeld.de/}{SGL}. Gemeinsam mit der Abteilung im Bewegungszentrum Langfort ist sie beim Deutschen Karate-Verband (\href{https://www.karate.de/}{DKV}), sowie dem Karate Dachverband Nordrhein-Westfalen (\href{https://www.karate.nrw/}{KDNW}) als offizieller Verein eingetragen. An beiden Standorten wird die Stilrichtung G\={o}j\={u}-Ry\={u} Karate-D\={o} gelehrt und praktiziert, es bestehen jedoch geringfügige Unterschiede in der jeweiligen Detailausprägung.
	\vspace{4pt}\newline
		Diese Prüfungsordnung orientiert sich an den offiziellen Vorgaben des \href{https://www.karate.de/}{DKV} und des G\={o}j\={u}-Ry\={u} Karate Deutschland (\href{https://www.karate-gkd.de/}{GKD}). Individuelle Abänderungen können zudem getroffen werden, sofern einzelne Vorgaben nicht umsetzbar sind. Den Vereinen ist es jedoch vorbehalten, das Prüfungsprogramm individuell anzupassen und in Ausnahmefällen zu verändern.
	\vspace{4pt}\newline
		In diesem Dokument wird die Prüfungsordnung der \href{https://www.sglangenfeld.de/}{SGL}-Karate Abteilung, für den Standort Reusrath, für Jugendliche und Erwachsene aufgeführt. Änderungen durch die Trainer und Prüfer sind jedoch jederzeit realisierbar und nicht individuell aufzuführen.
	\section{Allgemeine Voraussetzungen zur Prüfungsteilnahme}
		Um an einer Prüfung der Karate Abteilung am Standort Reusrath teilzunehmen sind unterschiedliche, individuelle Aspekte zu beachten. Grundsätzlich bietet die Abteilung alle fünf bis sieben Monate eine Gürtelprüfung an, wobei niemand gezwungen ist, an dieser teilzunehmen. Entweder werden Prüfer aus anderen Dojos\,/\,aus dem Verband eingeladen, oder die Trainer des Standorts Reusrath prüfen eigenständig. Aktuell sind Michael Gans und Lars Klein als Prüfer entsprechend vom \href{https://www.karate.de/}{DKV} lizenziert. Für alle Prüfungen gelten dabei die selben Grundsätze und Vorschriften, welche im Kommenden aufgeführt sind.
	\vspace{4pt}\newline
		Prüfungen im Karate-D\={o} dienen zum Einen dem Nachweis des Kenntnisstandes der Karate-Ka, zum anderen fungieren sie als Motivation, Herausforderung und Zielsetzung. Die einzelnen Prüfungsordnungen - spezifisch auf jeden Kyu-Grad übertragen - sind dabei an das vorausgesetzte, technische Niveau angepasst. Zum Bestehen einer Prüfung sind damit die technischen Kenntnisse elementar.
		\begin{center}
			\textbf{Niemand, der die Techniken und Vorgaben des kommenden Kyu-Grades nicht beherrscht oder umsetzen kann, kann an der Prüfung teilnehmen!}
		\end{center}
		Karate ist ein Prozess, welcher zwar anhand der Gürtelfarben ersichtlich wird, jedoch nur durch regelmäßiges und aufmerksames Training reift. Die Prüfungen symbolisieren folglich einen Fortschritt, der nicht innerhalb weniger Trainingsstunden zustande kommen kann.
	\vspace{4pt}\newline
		In unserem Verein werden demnach die individuellen Trainingsstunden jedes Mitgliedes dokumentiert (Trainingspass oder Liste), um vor Antritt der Prüfung einen Trainingsnachweis vorliegen zu haben. Grundsätzlich gilt dabei \textbf{immer}:
		\begin{center}
			\textbf{33\% der möglichen Trainingsstunden müssen absolviert worden sein.}
		\end{center}
		Wie viele Stunden dies letztlich sind, geben die Trainer etwa einen Monat vor der Prüfung bekannt. Wird dieser Wert nicht erreicht, muss auf eine Prüfung verzichtet werden und folglich an dem darauffolgenden Termin nicht teilgenommen werden. Zudem ist es den Trainern auch nach erfolgreich nachgewiesener Stundenzahl (33\% oder mehr) möglich, dem Karate-Ka von der Prüfung abzuraten oder diesen davon auszuschließen.
		\begin{center}
			\textbf{Einzig und Allein die Trainer sind es, welche die Karate-Ka zu einer Prüfung zulassen.\\(auch in Absprache externer Prüfer des \href{https://www.karate.de/}{DKV})}
		\end{center}
		Dabei werden sowohl die Trainingsstunden, als auch die technischen Voraussetzungen in Betracht gezogen. Obgleich den Trainern und Prüfern bewusst ist, dass eine Prüfung immer eine Motivation und Herausforderung für die Trainierenden darstellt, sind es stets die karatespezifischen Voraussetzungen welche erfüllt werden müssen. Ebenso ist es in der Karate Abteilung am Standort Reusrath keine Schande, nicht an einer Prüfung teilzunehmen, egal aus welchem Grund.
		\begin{center}
			\textbf{Wer sich nicht \frqq reif\flqq~genug für die nächste Prüfung fühlt, kann jederzeit aussetzen oder pausieren.}
		\end{center}
